% !Mode:: "TeX:UTF-8"

\theoremstyle{plain}
\theorembodyfont{\kai}
\theoremheaderfont{\kai\rmfamily}
\newtheorem{theorem}{\hei 定理}[chapter]
\newtheorem{axiom}{\hei 公理}[chapter]
\newtheorem{proposition}{\hei 性质}[chapter]
\newtheorem{lemma}{\hei 引理}[chapter]
\newtheorem{corollary}{\hei 推论}[chapter]

\theoremstyle{definition}
\theorembodyfont{\fs}
\theoremheaderfont{\kai\rmfamily}
\newtheorem{definition}{\hei 定义}[chapter]
\newtheorem{example}{\hei 例}[chapter]
\newtheorem{algo}{\hei 算法}[chapter]
\newtheorem{remark}{\hei 注解}[chapter]
%\newenvironment{proof}{\noindent{\kai 证明：}}{\hfill $ \square $ \vskip 4mm}
\newenvironment{proof}{}{\hfill $ \square $ \vskip 4mm}
\theoremsymbol{$\square$}

\allowdisplaybreaks[4]

\setlength{\parindent}{2em}

\arraycolsep=1.6pt

\renewcommand\contentsname{\hei 目~~~~录}

%\renewcommand\chaptername{\CJKprechaptername~\thechapter~\CJKchaptername}
\renewcommand{\chaptername}{第~\thechapter~章} % 自定义章节名称格式
%\setCJKmainfont{SimSun} % 设置中文字体

\setcounter{secnumdepth}{4} \setcounter{tocdepth}{2}

\titleformat{\chapter}{\center\fontsize{16pt}{16pt}\selectfont\hei}{第~\thechapter~章}{1em}{}
%%%%%%%%%%%%%%%%%%%%\titleformat{\chapter}{\center\xiaoer\hei}{\chaptertitlename}{0.5em}{}
\titlespacing*{\chapter}{0pt}{-5.5mm}{8mm}

%%%%%%%%%%%%%%%%%%%%%%%%%%%%%\titleformat{\section}{\xiaosan\hei}{\thesection}{0.5em}{}
\titleformat{\section}{\xiaosan\hei}{\thesection}{0.5em}{}
\titlespacing{\section}{0pt}{4.5mm}{4.5mm}
\titleformat{\subsection}{\sihao\hei}{\thesubsection}{0.5em}{}
\titlespacing{\subsection}{0pt}{4mm}{4mm}
\titleformat{\subsubsection}{\xiaosi\hei}{\thesubsubsection}{0.5em}{}
\titlespacing{\subsubsection}{0pt}{0pt}{0pt}
\titlecontents{chapter}[3.8em]{\hspace{-3.8em}\hei}{第~\thecontentslabel~章~~}{}{\titlerule*[4pt]{.}\contentspage}


\dottedcontents{section}[40pt]{}{22pt}{0.3pc}
\dottedcontents{subsection}[62pt]{}{32pt}{0.3pc}


\makeatletter
\renewcommand*\l@chapter{\@dottedtocline{0}{0em}{4em}}%控制英文目录： 细点\@dottedtocline  粗点\@dottedtoclinebold
\renewcommand*\l@section{\@dottedtocline{1}{1.5em}{1.8em}}
\renewcommand*\l@subsection{\@dottedtocline{2}{2.5em}{2.5em}}
%%页边距


% 定义页眉和页脚
\newcommand{\makeheadrule}{
\rule[7pt]{\textwidth}{0.75pt} \\[-23pt]
\rule{\textwidth}{2.25pt}}
\renewcommand{\headrule}{
    {\if@fancyplain\let\headrulewidth\plainheadrulewidth\fi
   \makeheadrule}}
\pagestyle{fancyplain}



%去掉章节标题中的数字
%%不要注销这一行，否则页眉会变成：“第1章1  绪论”样式
\renewcommand{\chaptermark}[1]{\markboth{第~\thechapter~章\quad #1}{}}

\fancyhf{}

%在book文件类别下,\leftmark自动存录各章之章名,\rightmark记录节标题
%% 页眉字号 工大要求 小五
%根据单双面打印设置不同的页眉；

\fancyhead[CO]{\song \xiaowu\leftmark}
\fancyhead[CE]{\song \xiaowu 数据分析与可视化 · 课程期末项目报告 }% === 课程项目适配：原“上海财经大学\cxuewei 学位论文” ===
\fancyfoot[C,C]{\xiaowu -~\thepage~-}


\renewcommand\frontmatter{\cleardoublepage
  \@mainmatterfalse
  \pagenumbering{Roman}}

% 设置行距和段落间垂直距离
%\renewcommand{\CJKglue}{\hskip 0.96pt plus 0.08\baselineskip} %加大字间距，使每行33个字
%\renewcommand{\captionfont}{\linespread{1.6}\normalsize}
%\renewcommand{\baselinestretch}{1}
%\setlength{\baselineskip}{20pt}
% 调整罗列环境的布局
\setitemize{leftmargin=3em,itemsep=0em,partopsep=0em,parsep=0em,topsep=-0em}
\setenumerate{leftmargin=3em,itemsep=0em,partopsep=0em,parsep=0em,topsep=0em}

\newcommand{\citeup}[1]{\textsuperscript{\cite{#1}}}

% 定制浮动图形和表格标题样式
\captionnamefont{\wuhao}
\captiontitlefont{\wuhao}
\captiondelim{~~}
\captionstyle{\centering}
\renewcommand{\subcapsize}{\wuhao}
\setlength{\abovecaptionskip}{0pt}
\setlength{\belowcaptionskip}{0pt}

% 自定义项目列表标签及格式 \begin{publist} 列表项 \end{publist}
\newcounter{pubctr} %自定义新计数器
\newenvironment{publist}{%%%%%定义新环境
\begin{list}{[\arabic{pubctr}]} %%标签格式
    {
     \usecounter{pubctr}
     \setlength{\leftmargin}{2.5em}     % 左边界 \leftmargin =\itemindent + \labelwidth + \labelsep
     \setlength{\itemindent}{0em}     % 标号缩进量
     \setlength{\labelsep}{1em}       % 标号和列表项之间的距离,默认0.5em
     \setlength{\rightmargin}{0em}    % 右边界
     \setlength{\topsep}{0ex}         % 列表到上下文的垂直距离
     \setlength{\parsep}{0ex}         % 段落间距
     \setlength{\itemsep}{0ex}        % 标签间距
     \setlength{\listparindent}{0pt} % 段落缩进量
    }}
{\end{list}}%%%%%

% 默认字体
\renewcommand\normalsize{
  \@setfontsize\normalsize{12pt}{12pt}
  \setlength\abovedisplayskip{4pt}
  \setlength\abovedisplayshortskip{4pt}
  \setlength\belowdisplayskip{\abovedisplayskip}
  \setlength\belowdisplayshortskip{\abovedisplayshortskip}
  \let\@listi\@listI}
\def\defaultfont{\renewcommand{\baselinestretch}{1.62}\normalsize\selectfont}
\predisplaypenalty=0  %公式之前可以换页，公式出现在页面顶部

% 封面、摘要、版权、致谢格式定义
\def\ctitle#1{\def\@ctitle{#1}}\def\@ctitle{}
\def\cdegree#1{\def\@cdegree{#1}}\def\@cdegree{}
\def\caffil#1{\def\@caffil{#1}}\def\@caffil{}
\def\csubject#1{\def\@csubject{#1}}\def\@csubject{}
\def\cIDnumber#1{\def\@cIDnumber{#1}}\def\@cIDnumber{}
\def\cauthor#1{\def\@cauthor{#1}}\def\@cauthor{}
\def\csupervisor#1{\def\@csupervisor{#1}}\def\@csupervisor{}
\def\cassosupervisor#1{\def\@cassosupervisor{{\hei 副 \hfill 导 \hfill 师} & #1\\}}\def\@cassosupervisor{}
\def\ccosupervisor#1{\def\@ccosupervisor{{\hei 联 \hfill 合\hfill 导 \hfill 师} & #1\\}}\def\@ccosupervisor{}
\def\cdate#1{\def\@cdate{#1}}\def\@cdate{}
\long\def\cabstract#1{\long\def\@cabstract{#1}}\long\def\@cabstract{}
\def\ckeywords#1{\def\@ckeywords{#1}}\def\@ckeywords{}
\def\authorization#1{\def\@authorization{#1}}\def\@authorization{}

\def\etitle#1{\def\@etitle{#1}}\def\@etitle{}
\def\edegree#1{\def\@edegree{#1}}\def\@edegree{}
\def\eaffil#1{\def\@eaffil{#1}}\def\@eaffil{}
\def\esubject#1{\def\@esubject{#1}}\def\@esubject{}
\def\eIDnumber#1{\def\@eIDnumber{#1}}\def\@eIDnumber{}
\def\eauthor#1{\def\@eauthor{#1}}\def\@eauthor{}
\def\esupervisor#1{\def\@esupervisor{#1}}\def\@esupervisor{}
\def\eassosupervisor#1{\def\@eassosupervisor{\textbf{Associate Supervisor:} & #1\\}}\def\@eassosupervisor{}
\def\ecosupervisor#1{\def\@ecosupervisor{\textbf{Co Supervisor:} & #1\\}}\def\@ecosupervisor{}
\def\edate#1{\def\@edate{#1}}\def\@edate{}
\long\def\eabstract#1{\long\def\@eabstract{#1}}\long\def\@eabstract{}
\long\def\NotationList#1{\long\def\@NotationList{#1}}\long\def\@NotationList{}
\def\ekeywords#1{\def\@ekeywords{#1}}\def\@ekeywords{}
\def\natclassifiedindex#1{\def\@natclassifiedindex{#1}}\def\@natclassifiedindex{}
\def\internatclassifiedindex#1{\def\@internatclassifiedindex{#1}}\def\@internatclassifiedindex{}
\def\statesecrets#1{\def\@statesecrets{#1}}\def\@statesecrets{}

% 定义封面
\def\makecover{
\pdfbookmark[0]{封面}{Cover}
\begin{titlepage}
\begin{center}
\includegraphics[width=0.7\textwidth]{setup/shufe.png}% === 适配:实际文件为png ===
%\vspace{2em}

     \parbox[t][2.4cm][t]{\textwidth}{\zihao{-1}
    \begin{center} {\hei  课程期末项目报告 }\end{center} }% === 课程项目适配：原“毕业论文” ===

    \parbox[t][5cm][t]{\textwidth}{\erhao
    \begin{center} {\hei \@ctitle}\end{center} }

	\parbox[t][8.8cm][b]{\textwidth}
     {\sihao
    \begin{center} \renewcommand{\arraystretch}{1.7} \songti\sihao
    \begin{tabular}{r@{：}c}
     {\hei 姓\hfill 名} & \@cauthor\\\cline{2-2}
     {\hei 学\hfill 号} & \@cIDnumber\\\cline{2-2}
    {\hei 学 \hfill 院}& \qquad 统计与管理学院\qquad \\\cline{2-2}
    {\hei 专\hfill 业} & \@csubject\\\cline{2-2}
     {\hei 指导老师} &\@csupervisor\\\cline{2-2}
     \end{tabular} \renewcommand{\arraystretch}{1}
    \end{center} }

    \vspace{3em}

   {\songti\sihao 定稿日期：~\chinesethesistime}
\end{center}
\end{titlepage}
%%%%%%%%%%%%%%%%%%%%%%%%%%%%%%%%%%%%%%%%%%%%%%%%%%%%%%%%%%%%%%%%%%%
% === 课程项目适配：注释掉学位论文诚信声明页，如需恢复取消注释即可 ===
%\pdfbookmark[0]{声明}{Statement}
%\begin{titlepage}
%\vspace{5em}
%\begin{center}{\hei\yihao 声\hspace{3cm}明}\end{center}
%
%\vspace{5em}
%
%
%{\kai\erhao
%本人郑重声明所呈交的论文是我个人在指导老师的指导下进行的研究工作及取得的研究成果，不存在任何剽窃、抄袭他人学术成果的现象。我同意（\hspace{2cm}   ）/不同意（  \hspace{2cm} ）本论文作为学校的信息资料使用。
%
%
%\vspace{5em}
%
%
%\hfill 论文作者（签名）\hspace{4em}
%
%\vspace{3em}
%
%
%\hfill 年\hspace{1.5em}月\hspace{1.5em}日
%}
%
%\thispagestyle{empty}
%\end{titlepage}

%%%%%%%%%%%%%%%%%%%   Abstract and keywords  %%%%%%%%%%%%%%%%%%%%%%%
\clearpage

\BiAppendixChapter{摘\quad 要}{Abstract (In Chinese)}

\setcounter{page}{1}
\song\defaultfont
\@cabstract

\vspace{\baselineskip}

\hangafter=1\hangindent=52.3pt\noindent
{\hei 关键词：} \@ckeywords

%%%%%%%%%%%%%%%%%%%   English Abstract  %%%%%%%%%%%%%%%%%%%%%%%%%%%%%%
\clearpage

\phantomsection
\markboth{Abstract}{Abstract}
\addcontentsline{toc}{chapter}{\xiaosi ABSTRACT}
\addcontentsline{toe}{chapter}{\bfseries \xiaosi Abstract (In English)}  \chapter*{\textbf{Abstract}}
\@eabstract
\vspace{\baselineskip}

\hangafter=1\hangindent=60pt\noindent
{\textbf{Keywords:}}  \@ekeywords
\clearpage
 \phantom{s}
 \thispagestyle{empty}
}
%%%%%%%%%%%%%%%%%%%%%%%%%%%%%%%%%%%%%%%%%%%%%%%%%%%%%%%%%%%%%%%
% 英文目录格式
\def\@dotsep{0.75}           % 定义英文目录的点间距
\setlength\leftmargini {0pt}
\setlength\leftmarginii {0pt}
\setlength\leftmarginiii {0pt}
\setlength\leftmarginiv {0pt}
\setlength\leftmarginv {0pt}
\setlength\leftmarginvi {0pt}

\def\engcontentsname{\bfseries Contents}
\newcommand\tableofengcontents{
   \pdfbookmark[0]{Contents}{econtent}
     \@restonecolfalse
   \chapter*{\engcontentsname  %chapter*上移一行，避免在toc中出现。
       \@mkboth{%
          \engcontentsname}{\engcontentsname}}
   \@starttoc{toe}%
   \if@restonecol\twocolumn\fi
   }

\urlstyle{same}  %论文中引用的网址的字体默认与正文中字体不一致，这里修正为一致的。

\renewcommand\endtable{\vspace{-4mm}\end@float}

%%%定义附录

\makeatother
